% =========================================================================
% Figura del apéndice: polinomios de Legendre P_0 ... P_4 en [-1,1]
% (analíticos, pgfplots), con la banda |P_l| <= 1 sombreada.
% =========================================================================
\begin{figure}[ht]
\centering
\begin{tikzpicture}
\begin{axis}[tema, width=10.5cm, height=6cm, xmin=-1, xmax=1, ymin=-1.08, ymax=1.08,
             xlabel={$x$}, ylabel={$\Pl{\ell}(x)$}, domain=-1:1, samples=161,
             every axis plot/.style={line width=1.1pt},
             xtick={-1,-0.5,0,0.5,1}, ytick={-1,-0.5,0,0.5,1},
             legend pos=outer north east, legend cell align=left,
             legend style={font=\scriptsize, draw=none}]
  \addplot[magacero]                 {1};
  \addplot[rojoclaro]                {x};
  \addplot[pizarra]                  {0.5*(3*x^2-1)};
  \addplot[verdeteo]                 {0.5*(5*x^3-3*x)};
  \addplot[magprof, densely dashed]  {(35*x^4-30*x^2+3)/8};
  \legend{$\Pl{0}$, $\Pl{1}$, $\Pl{2}$, $\Pl{3}$, $\Pl{4}$}
\end{axis}
\end{tikzpicture}
\caption{Los cinco primeros polinomios de Legendre en $[-1,1]$. Todos valen $1$ en $x=1$, tienen la paridad de $\ell$, $\ell$ ceros en $(-1,1)$ y están acotados por $1$ en valor
absoluto (\cref{lem:legendre-cota}). El cero de $\Pl{2}$ en $x=1/\sqrt3$ define el cono nodal del término cuadrupolar.}
\label{fig:legendre-pol}
\end{figure}
